% Slides for advisor meeting: core of the model.
% Compile with pdflatex + bibtex; needs references.bib (the paper's) in the same folder.
\documentclass[11pt,aspectratio=169]{beamer}

\usepackage[utf8]{inputenc}
\usepackage{amsmath,amssymb,mathtools}
\usepackage{booktabs}
\usepackage{graphicx}
\usepackage{natbib}
\usepackage{tikz}
\usetikzlibrary{arrows.meta}

% Look: same as the paper (Computer Modern serif, black text, blue links)
\usefonttheme{serif}
\setbeamertemplate{navigation symbols}{}
\setbeamercolor{structure}{fg=black}
\setbeamercolor{alerted text}{fg=blue!75!black}
\setbeamerfont{frametitle}{series=\bfseries}
\setbeamerfont{framesubtitle}{size=\footnotesize,shape=\itshape}
\setbeamertemplate{itemize items}[circle]
\setbeamertemplate{bibliography item}{}
\setbeamertemplate{frametitle continuation}{}
\setbeamerfont{block title}{series=\bfseries}
\hypersetup{
    colorlinks=true,
    linkcolor=blue,
    filecolor=magenta,
    urlcolor=cyan,
    pdftitle={The Political Economy of State Formation after the Introduction of Agriculture},
    }

% Footer as in the paper: rule, page number on the right
\setbeamerfont{footline}{size=\scriptsize}
\setbeamertemplate{footline}{%
  \hbox to \paperwidth{\hskip1.5em\vbox{\hsize=\dimexpr\paperwidth-3em\relax
    \hrule height 0.6pt\vskip0.6ex
    \hbox to \hsize{\hfill\usebeamerfont{footline}Page \insertframenumber}%
    \vskip1.2ex}\hss}}

\newtheorem{proposition}[theorem]{Proposition}

\title{The Political Economy of State Formation\\after the Introduction of Agriculture}
\author{Alberto Bisin \and Dionysios Papadakis}
\date{\today}

\begin{document}

\begin{frame}[plain,noframenumbering]
\titlepage
\end{frame}

%----------------------------------------------------------------------
\begin{frame}{What the model does}
\begin{itemize}
\item A \textbf{priesthood} builds and runs the central organisation of the polity: storage,
redistribution, the calendar, ritual. This is \textbf{state capacity} $s$.
\item Capacity is costly and the priests are paid \emph{only} through taxes. There is no
commitment, so too little is built.
\item Farmers can give the priests more \textbf{political weight} $\beta$ $\Rightarrow$ higher
taxes $\Rightarrow$ higher return to building capacity.\\
Farmers delegate power \alert{because they value the capacity it induces}.
\item Only \textbf{religious} farmers value the temple's services. Religion is transmitted as in
\citet{bisin2001}, so it spreads when the temple delivers.
\item \textbf{Legibility} $\omega$ limits how much of the harvest can be taxed at all.
\item Institutions are political weights updated as in \citet{bisin2024joint}. Culture,
institutions and capacity co-evolve.
\end{itemize}
\end{frame}

%----------------------------------------------------------------------
\begin{frame}{Patterns we want the model to speak to}
\begin{enumerate}
\item \textbf{Religious authority before administration.} Monumental Ubaid temples (Eridu) with no
writing; writing appears in Late Uruk, as temple accounts. Same order in the Andes (Norte Chico).
\item \textbf{Religious authority is common, states are rare.} Hierarchy rises with the
appropriability of the staple crop \citep{mayshar2022}, at the chiefdom$\to$state margin but not
at tribe$\to$chiefdom \citep{cook2026}.
\item \textbf{Long stasis, fast transitions.} About six thousand years from the first storage to the
first states in the Fertile Crescent; the transition itself takes a few centuries.
\item \textbf{The early state's services are religious.} The temple is the household of the god;
taxation and redistribution are the management of divine property.
\end{enumerate}
\end{frame}

%----------------------------------------------------------------------
\section{Model}

\begin{frame}{Environment}
Discrete time. A unit mass of farmers $F$ and a priesthood $P$ (a single agent) that runs the
temple, which is also the central organisation of the polity.

\medskip
\begin{columns}[T]
\column{0.44\textwidth}
\textbf{State} $(s_t,\beta_t,q_t)$:
\begin{itemize}
\item $s_t\ge0$: state capacity;
\item $\beta_t\in[0,1]$: priests' weight;
\item $q_t\in[0,1]$: religious share.
\end{itemize}
\column{0.52\textwidth}
\centering\small
\begin{tabular}{l|cc}
 & $\beta$ low & $\beta$ high \\
\hline
$s$ low & acephalous & ritual authority\\[3pt]
$s$ high & limited delegation & temple state
\end{tabular}
\end{columns}

\medskip
\textbf{Within period $t$}, given $(s_{t-1},\beta_t,q_t)$:
\begin{enumerate}
\item the priests invest $i_t$ and the political process sets the tax $\tau_t$, simultaneously;
\item religious farmers decide whether to join the cult;
\item culture updates $q_{t+1}$ and institutions update $\beta_{t+1}$.
\end{enumerate}
\end{frame}

%----------------------------------------------------------------------
\begin{frame}{Farmers}
Endowment $Y$, consumption $c_t=(1-z_t)Y$, where $z_t$ is the effective tax rate. Two cultural
types $i\in\{r,n\}$. Joining the cult, $a\in\{0,1\}$, costs $e>0$:
\[
U^i_t=\log c_t+a\bigl(\theta_i\,v(G_t)-e\bigr),\qquad \theta_r>\theta_n=0,
\]
$G_t$ are the temple's services; $v$ increasing, strictly concave, $v(0)=0$, $v'(0)<\infty$
(numerics: $v(G)=\log(1+G)$).
\begin{itemize}
\item Non-religious farmers never participate. Religious farmers participate iff
$G_t\ge\underline s\equiv v^{-1}(e/\theta_r)$.
\item Both types consume the same, so aggregate farmer welfare is
\[
\log c_t+q_t\,w(G_t),\qquad w(G)\equiv\max\{\theta_rv(G)-e,\,0\}.
\]
\item \alert{Legitimacy} $\equiv q_t\theta_r$: the weight society puts on the temple's services.
It enters only preferences, not reach, the tax base or any cost.
\end{itemize}
\end{frame}

%----------------------------------------------------------------------
\begin{frame}{Priests and capacity}
The priests invest $i_t\ge0$ (available within the period):
\[
s_t=(1-\delta)s_{t-1}+i_t,\qquad \delta\in(0,1),
\]
and consume revenue net of the cost of investment:
\[
U^P_t=T_t-\frac{\psi}{2}\,i_t^2,\qquad \psi>0 .
\]
\begin{itemize}
\item Services equal the capacity stock, $G_t=s_t$: the granaries, herds, officials and records
that provide the services are the same ones that collect taxes.
\item The priests are the only builders, and they are paid only out of $T_t$.
\end{itemize}
\end{frame}

%----------------------------------------------------------------------
\begin{frame}{Reach and taxation}
Share of the endowment within fiscal reach:
\[
\rho(s,\omega)=\omega\,r(s),\qquad r(s)=\frac{\kappa+s}{1+s},\qquad
\omega\in(0,1],\ \kappa\in(0,1).
\]
\begin{itemize}
\item $\omega$: \alert{legibility}, how much of the harvest can be assessed at all (cereals vs.\
roots and tubers, \citealp{mayshar2017,mayshar2022}).
\item $\kappa$: reach by presence alone, without records or officials.
\item All that matters: $r$ increasing, strictly concave, $r(0)=\kappa>0$; $\omega$ scales reach
proportionally.
\end{itemize}
Statutory rate $\tau_t\in[0,1]$ on activity within reach:
\[
z_t=\tau_t\,\rho(s_t,\omega),\qquad T_t=z_tY .
\]
Two limits on taxes: $\log c\to-\infty$ as $z\to1$, and $z\le\rho(s,\omega)<\omega\le1$.
\end{frame}

%----------------------------------------------------------------------
\begin{frame}{Within-period game: no commitment}
Given $(s_{t-1},\beta_t,q_t)$, \emph{simultaneously}:
\begin{itemize}
\item the political process chooses $\tau$, taking $i_t$ as given, to maximise
\[
\beta_t\Bigl[\tau\rho(s_t,\omega)Y-\tfrac{\psi}{2}i_t^2\Bigr]
+(1-\beta_t)\Bigl[\log\bigl((1-\tau\rho(s_t,\omega))Y\bigr)+q_t\,w(s_t)\Bigr];
\]
\item the priests choose $i_t$, taking $\tau_t$ as given, to maximise
$\tau_t\,\rho(s_t,\omega)\,Y-\tfrac{\psi}{2}i_t^2$.
\end{itemize}
\medskip
\textbf{The friction} \citep[as in][]{bisin2024joint}: the political process cannot commit to a tax
before capacity is built, and the priests cannot contract with farmers over investment.
\begin{itemize}
\item Farmers value capacity, but only the priests pay for it, and they are paid only through
taxes.
\item The political process trades off the priests' extraction against farmers' consumption; it
does not internalise the value of capacity.
\end{itemize}
\end{frame}

%----------------------------------------------------------------------
\begin{frame}{Societal equilibrium}
\begin{lemma}[Societal equilibrium]
For every $(s_{t-1},\beta_t)$ the equilibrium exists and is unique, with
\[
z_t=\min\bigl\{z(\beta_t),\ \rho(s_t,\omega)\bigr\},\qquad
z(\beta)\equiv\max\Bigl\{0,\,1-\frac{1-\beta}{\beta Y}\Bigr\}.
\]
\begin{enumerate}[(i)]
\item \emph{Interior}: $\tau_t<1$, $z_t=z(\beta_t)$, $\psi i_t=z(\beta_t)\,Y\,r'(s_t)/r(s_t)$:
no $\omega$.
\item \emph{Reach-constrained}: $\tau_t=1$, $z_t=\rho(s_t,\omega)$, $\psi i_t=\omega\,r'(s_t)\,Y$:
no $\beta_t$.
\end{enumerate}
Regime (ii) iff $\beta_t\ge\hat\beta(s_{t-1})\equiv\bigl[1+Y\bigl(1-\rho(s^R_t,\omega)\bigr)\bigr]^{-1}$,
where $s^R_t$ is the capacity the priests build at $\tau=1$.
\end{lemma}
\begin{itemize}
\item In both regimes $\psi i_t=T_t\,r'(s_t)/r(s_t)$: whatever raises the tax base raises
investment.
\item No tax for $\beta\le\underline\beta\equiv1/(1+Y)$. Weight above $\hat\beta$ is
\emph{nominal}: effective power is $\beta^{\rm eff}_t=\min\{\beta_t,\hat\beta(s_{t-1})\}$.
\end{itemize}
\end{frame}

%----------------------------------------------------------------------
\begin{frame}{Legibility and investment}
\begin{lemma}[Legibility and investment]
Fix $(s_{t-1},\beta_t)$ with $\beta_t>\underline\beta$. There is a threshold $\hat\omega$,
increasing in $\beta_t$, such that
\begin{enumerate}[(i)]
\item if $\omega<\hat\omega$, legibility binds: $\tau_t=1$ and
$\dfrac{\partial i_t}{\partial\omega}=\dfrac{r'(s_t)\,Y}{\psi-\omega\,r''(s_t)\,Y}>0$;
\item if $\omega>\hat\omega$, legibility does not bind: $\tau_t<1$ and $i_t$ does not depend on
$\omega$.
\end{enumerate}
Investment is continuous, hence non-decreasing, in $\omega$.
\end{lemma}
\begin{itemize}
\item Legibility is a \emph{ceiling}. When it binds, raising it raises the priests' return to
capacity.
\item When it does not bind, the political process offsets higher $\omega$ with a lower $\tau$:
$z=\tau\rho$ and the return to capacity are unchanged.
\item $\hat\omega$ rises with $\beta_t$: legibility matters over a wider range the more power the
priests hold.
\end{itemize}
\end{frame}

%----------------------------------------------------------------------
\begin{frame}{Institutional change}
Political weights are updated between periods by the myopic rule of \citet{bisin2024joint}:
\[
\beta_{t+1}\in\arg\max_{b\in[0,1]}\ \Bigl\{\beta_t\,V_P(b)+(1-\beta_t)\,V_F(b)\Bigr\},
\]
\[
V_P(b)=z(b)Y-\tfrac{\psi}{2}\,i(b)^2,\qquad
V_F(b)=\log\bigl((1-z(b))Y\bigr)+q_t\,w\bigl(s(b)\bigr),
\]
the indirect utilities in the societal equilibrium at weight $b$, starting from $s_{t-1}$.
\begin{itemize}
\item Current weights evaluate alternative weights: institutions internalise effects that the
within-period equilibrium does not.
\item The objective is not concave in $b$ (flat below $\underline\beta$ and above $\hat\beta$,
kink where $s(b)$ crosses $\underline s$), so the update compares \emph{all} weights, not only
nearby ones.
\end{itemize}
\end{frame}

%----------------------------------------------------------------------
\begin{frame}{The delegation margin: what changes with $b$}
Fix $(s_{t-1},q_t)$. $\bigl(z(b),i(b),s(b)\bigr)$ is period $t$'s societal equilibrium \emph{if}
the priests' weight were $b$. Interior regime, $\underline\beta<b<\hat\beta(s_{t-1})$:
\[
\underbrace{z(b)=1-\frac{1-b}{bY}}_{\text{tax rule}},\qquad
\underbrace{\psi\,i(b)=z(b)\,Y\,\frac{r'(s(b))}{r(s(b))}}_{\text{priests' FOC at the equilibrium tax}},\qquad
s(b)=(1-\delta)s_{t-1}+i(b).
\]
\begin{itemize}
\item $i(\cdot)$ is the \textbf{equilibrium investment map}; $i_t\equiv i(\beta_t)$ and
$s_t\equiv s(\beta_t)$ are this period's values. Its slope:
\[
i'(b)=\frac{z'(b)\,Y\,(r'/r)(s(b))}{\psi-z(b)\,Y\,(r'/r)'(s(b))}>0
\qquad\bigl(z'(b)=1/(b^2Y)>0,\ (r'/r)'<0\bigr).
\]
Channel: $b\uparrow\Rightarrow z\uparrow\Rightarrow$ the priests' marginal return
$\tau\rho_sY=zY\,r'/r$ rises (they take $\tau$ as given) $\Rightarrow i\uparrow$.
\item Revenue is $z(b)Y$ for any $s$, since $\tau=z(b)/\rho(s)$ adjusts. So
$V_P(b)=z(b)Y-\tfrac\psi2 i(b)^2$: induced investment is a pure cost to the priests.
\end{itemize}
\end{frame}

%----------------------------------------------------------------------
\begin{frame}{The delegation margin: two effects}
Updating objective $W(b)\equiv\beta_tV_P(b)+(1-\beta_t)V_F(b)$: the weights are fixed at the
current $\beta_t$; only the allocation moves with $b$. In the interior regime, with religious
farmers participating ($s(b)>\underline s$),
\[
W'(b)=
\underbrace{z'(b)\Bigl[\beta_t\,Y-\frac{1-\beta_t}{1-z(b)}\Bigr]}_{\text{transfer effect}}
+\underbrace{i'(b)\Bigl[(1-\beta_t)\,q_t\theta_r\,v'(s(b))-\beta_t\,\psi\,i(b)\Bigr]}_{\text{capacity effect}} .
\]
\begin{itemize}
\item \textbf{Transfer effect.} More weight raises the tax by $z'(b)$. Priests gain $Y$ per unit
of tax, farmers lose $1/(1-z)$ in consumption utility. At $b=\beta_t$ the tax rule equates the
two, so this term is \emph{zero}. It is negative for $b>\beta_t$: second order for small
delegations, first order for large ones.
\item \textbf{Capacity effect.} More weight induces $i'(b)$ more investment.
\begin{itemize}
\item Farmers gain services worth $q_t\theta_r\,v'(s)$ per unit: only religious farmers, and
only if they participate.
\item Priests pay the marginal cost $\psi\,i$ and get no extra revenue.
\end{itemize}
\end{itemize}
\end{frame}

%----------------------------------------------------------------------
\begin{frame}{The delegation margin}
\begin{lemma}[Delegation margin]
In the interior regime, with religious farmers participating,
\[
W'(\beta_t)= i'(\beta_t)\Bigl[(1-\beta_t)\,q_t\theta_r\,v'(s_t)-\beta_t\,\psi\,i_t\Bigr],
\qquad i'(\beta_t)>0 .
\]
Without participation it equals $-i'(\beta_t)\,\beta_t\,\psi\,i_t<0$.
\end{lemma}
\begin{itemize}
\item It balances out two effects of the extra investment $i'(\beta_t)$ that more power induces:
farmers value additional elite power to the extent that they value the public good,
$q_t\theta_rv'(s_t)$, with weight $1-\beta_t$; priests lose from it to the extent that they pay
for building it, $\psi i_t$, with weight $\beta_t$.
\[
\boxed{\ W'(\beta_t)>0\iff(1-\beta_t)\,q_t\theta_r\,v'(s_t)>\beta_t\,\psi\,i_t\ }
\]
\item \emph{Local} statement: the update is global, so large moves can still win. An interior
steady state needs equality.
\item With no believers $W'(\beta_t)<0$: power falls until the priests cannot tax,
$\beta\le\underline\beta$.
\end{itemize}
\end{frame}

%----------------------------------------------------------------------
\begin{frame}{No commitment is the friction}
\begin{proposition}[Commitment benchmark]
Suppose the political process chooses $\tau_t$ before the priests invest, anticipating their best
response. Then for every $\beta_t$ the updating objective is maximised at $b=\beta_t$: political
weights never change, and capacity, if any, is built at the prevailing weights.
\end{proposition}
\begin{itemize}
\item Under commitment, the allocation chosen at any weight $b$ is also available to the process at
weight $\beta_t$, which picks the best one for $\beta_t$. No $b$ can beat $\beta_t$.
\item So delegation is a \alert{substitute for commitment}: farmers cannot promise the priests a
return on capacity, so they give them the power to take it.
\end{itemize}
\end{frame}

%----------------------------------------------------------------------
\begin{frame}{Cultural change}
Socialisation as in \citet{bisin2001}: a child of a type-$i$ parent is socialised by the parent
with probability $d^i$, and otherwise adopts the trait of a randomly drawn adult. Parents are
imperfectly empathetic and evaluate the child's life at current services.

\medskip
Only participation differs across types, so
\[
\Delta V^r_t=\theta_r v(s_t)-e,\qquad \Delta V^n_t=e
\qquad(\text{if } s_t>\underline s;\ \text{both zero otherwise}).
\]
With quadratic socialisation costs, $d^i=\lambda(1-q^i)\Delta V^i$ and
\[
q_{t+1}-q_t=\lambda\,q_t(1-q_t)\bigl[(1-q_t)\Delta V^r_t-q_t\Delta V^n_t\bigr].
\]
\begin{itemize}
\item Religious parents want observance because they value the cult's services; non-religious
parents see observance as a pure cost.
\end{itemize}
\end{frame}

%----------------------------------------------------------------------
\begin{frame}{Religion follows the temple}
\begin{lemma}[Religion follows the temple]
If $s_t>\underline s$, the cultural dynamics have the unique interior attractor
\[
q^*(s_t)=1-\frac{e}{\theta_r\,v(s_t)},
\]
which is increasing in $s_t$. If $s_t\le\underline s$, neither type socialises and $q_{t+1}=q_t$:
religion is frozen, or latent.
\end{lemma}
\begin{itemize}
\item The \emph{fixed} observance cost $e$ is what makes religiosity respond to capacity. If the
cost were proportional to $v(G)$, the two motives would be proportional and $q^*$ would not
depend on $s$.
\item Simplifications: parents evaluate at current services, not the services the child will face;
religion has no effect on production or the tax base.
\end{itemize}
\end{frame}

%----------------------------------------------------------------------
\begin{frame}{The mechanism}
\centering
\begin{tikzpicture}[
  font=\small,
  state/.style={draw, thick, rounded corners=3pt, fill=black!7, align=center,
                minimum width=3cm, minimum height=0.9cm},
  lab/.style={align=center, font=\footnotesize},
  link/.style={-{Stealth[length=2.5mm]}, thick},
]
\node[state] (b) at (0,3.2)    {Political weight $\beta$};
\node[state] (s) at (-4,0) {Capacity $s$};
\node[state] (q) at (4,0)  {Religious share $q$};

\draw[link] (b.west) to[out=180,in=90] coordinate[pos=0.45] (m1) (s.north);
\draw[link] (s.east) -- node[lab, below=2pt] {religion follows the temple} (q.west);
\draw[link] (q.north) to[out=90,in=0] (b.east);

\node[lab] at (-1.9,1.5) {more power,\\ more investment};
\node[lab] at (1.9,1.5)  {farmers delegate\\ for capacity};

\node[lab, text=black!60] (om) at (-5.6,2.6) {legibility $\omega$\\ caps this link};
\draw[-{Bar[width=2mm]}, black!50] (om.east) to[out=0,in=160] (m1);
\end{tikzpicture}

\bigskip
\raggedright
\begin{itemize}
\item $\beta\to s$: societal equilibrium ($\psi i=T\,r'/r$), capped by $\omega$ at full reach.
\item $s\to q$: religion follows the temple, $q^*(s)$ increasing.
\item $q\to\beta$: delegation margin, farmers value capacity at $q\theta_rv'(s)$.
\end{itemize}
\end{frame}

%----------------------------------------------------------------------
\section{Results}

% Results slides: smaller frame titles so the long titles fit on one line
\setbeamerfont{frametitle}{size=\large,series=\bfseries}

\begin{frame}{Some preliminary results (Claude) - State Formation}
Start from an \emph{acephalous} polity: no inherited capacity, and priests too weak to tax.
\begin{itemize}
\item \textbf{The acephalous polity is a trap.} No tax, no revenue, no investment. A small
delegation builds too little capacity for anyone to join the cult, so it is a pure transfer to the
priests, which the farmers reject. Religion stays frozen.
\item \textbf{Leaving requires a big push.} Only a delegation that lifts capacity above the
participation threshold in one period can pay, and only if enough farmers are (latently)
religious: there is a threshold religious share.
\item \textbf{What lowers the threshold.} Inherited authority of the priests, a higher value of the
temple's services, a lower observance cost. In simulations also cheaper capacity, and inherited
capacity: a temple built earlier substitutes for believers.
\item \textbf{What raises it} (simulations). Higher base reach $\kappa$: when much can be taken
without administration, a given tax induces less building, so the transfer buys less of what
farmers value.
\end{itemize}
\end{frame}

%----------------------------------------------------------------------
\begin{frame}{Some preliminary results (Claude) - State Formation}
\begin{itemize}
\item \textbf{Legibility matters only through a floor.} Below a minimum legibility no state forms,
whatever the religious share. Above it, more legibility helps only while it limits how large a
delegation can be; beyond that it is irrelevant for formation.
\item \textbf{Forming is harder than keeping.} Forming needs capacity above the participation
threshold within one period; keeping a state needs only replacing depreciation. A state can be
sustainable but unreachable, so both outcomes can be long-run, and history decides.
\item \textbf{Long stasis, then a fast transition.} The acephalous polity is absorbing, so formation
waits for a slowly moving parameter (e.g.\ the cost of capacity falling as know-how accumulates)
to push the threshold below the religious share. Then capacity rises within a few periods. The
model explains the suddenness, not the timing.
\item \textbf{Sequencing} (simulations). Power moves first, capacity follows, and religiosity
adjusts last and most slowly.
\end{itemize}
\end{frame}

%----------------------------------------------------------------------
\begin{frame}{Some preliminary results (Claude) - Steady State Comp. Statics}
Comparative statics of the \emph{long-run state}: the stable steady state with positive capacity,
where investment just replaces depreciation, the religious share is at its attractor, and the
priests' weight no longer changes.
\begin{itemize}
\item \textbf{Size is set by the scarcer resource.} Long-run capacity is the smaller of two levels:
\begin{itemize}
\item a \emph{legitimacy} level $s^H$: the largest capacity at which the farmers' valuation of more
capacity still justifies the extra power needed to induce it;
\item a \emph{legibility} level $s^R(\omega)$: the capacity the priests maintain when they already
tax everything within reach, so more power buys nothing.
\end{itemize}
The state is \emph{legitimacy-limited} or \emph{legibility-limited}. Improving the scarce
resource raises capacity; improving the other has no effect.
\item \textbf{Power, tax and religion follow capacity.} The long-run tax is what the priests need
to maintain that capacity, power is the weight that delivers that tax, and religiosity is what a
temple of that size sustains. All three rise with long-run capacity.
\end{itemize}
\end{frame}

%----------------------------------------------------------------------
\begin{frame}{Some preliminary results (Claude) - Steady State Comp. Statics}
\begin{columns}[T]
\column{0.40\textwidth}
\centering
\begin{tabular}{lcc}
\toprule
 & $s^H$ & $s^R(\omega)$\\
\midrule
$\theta_r$ (valuation) & $+$ & $0$\\
$e$ (observance) & $-$ & $0$\\
$\omega$ (legibility) & $0$ & $+$\\
$Y$ (endowment) & $+$ & $+$\\
$\psi,\ \delta,\ \kappa$ & $-$ & $-$\\
\bottomrule
\end{tabular}

\smallskip
\raggedright\footnotesize
Effect of each parameter on the two levels; long-run capacity moves with whichever one binds.
\column{0.56\textwidth}
\begin{itemize}
\item Legibility raises capacity, the tax and power only if the state is legibility-limited;
$\theta_r$ and $e$ matter only if it is legitimacy-limited.
\item Religiosity rises with $\theta_r$ and falls with $e$ always. It rises with legibility only if
legibility-limited: the temple is larger, belief is not directly affected.
\end{itemize}
\end{columns}
\medskip
\begin{itemize}
\item \textbf{Base reach} $\kappa$ lowers capacity but, at the reach frontier, raises the tax and
power: more power over a smaller apparatus (the ritual-authority cell).
\item \textbf{Writing}, read as a rise in legibility, pays only for a polity that already taxes
everything within reach, and more so the larger its temple.
\item \textbf{Inherited authority persists.} Priests with weight above the kink keep it and maintain
frontier capacity even when legitimacy alone would support less.
\end{itemize}
\end{frame}

% Back to the default frame-title size
\setbeamerfont{frametitle}{size=\Large,series=\bfseries}

%----------------------------------------------------------------------
\begin{frame}{What is doing the work, and open issues}
\begin{columns}[T]
\column{0.48\textwidth}
\textbf{Assumptions that matter}
\begin{itemize}
\item No commitment within the period; with commitment weights never move.
\item Priests' payoff linear in revenue. With $U^P=\frac{T^{1-\sigma}}{1-\sigma}-\frac\psi2 i^2$,
$\psi i=T^{1-\sigma}r'/r$: the channel vanishes at $\sigma=1$.
\item Fixed observance cost $e$: makes $q^*$ depend on $s$ and creates the participation
non-convexity behind the trap.
\item $G=s$: one stock delivers both the services and the fiscal reach.
\end{itemize}
\column{0.48\textwidth}
\textbf{Open}
\begin{itemize}
\item Single-peakedness of $\Omega$ is only numerical.
\item The formation result covers the first delegation; that the polity then reaches $S_\infty$ is
numerical.
\item Steady states with nominal power ($\beta>\hat\beta$), the tie-break, and collapse.
\item Mapping to \citet{cook2026}: is ``formation vs.\ size'' the right analogue of
``tribe$\to$chiefdom vs.\ chiefdom$\to$state''?
\item No neighbours, exit or conquest.
\end{itemize}
\end{columns}
\end{frame}

%----------------------------------------------------------------------
\begin{frame}[allowframebreaks]{References}
\footnotesize
\bibliographystyle{apsr}
\bibliography{references}
\end{frame}

\end{document}
